\section{Round-five positive closure: recut Hamiltonian suspension and typed response}
\label{sec:r5-a1}

This section replaces the incompatible incoming--outgoing mapping-torus
identification by a single cut symplectic return section.  It also replaces the
non-descending piecewise constant one-form by a global collar-supported work
form and defines every operator used in the response calculus.

\subsection{The common cut section and the recut map}

Let
\[
 0=s_0(a)<s_1(a)<\cdots<s_4(a)=1,
 \qquad p_i(a)=s_i(a)-s_{i-1}(a)>0,
\]
and put \(I_i(a)=(s_{i-1}(a),s_i(a))\).  Let
\(0=t_0(a)<\cdots<t_4(a)=1\) have the same increments and put
\(J_i(a)=(t_{i-1}(a),t_i(a))\).  On the incoming port \(I_i(a)\times(0,1)\)
set
\[
 R_{a,i}(q,p)=
 \left(\frac{q-s_{i-1}(a)}{p_i(a)},
       t_{i-1}(a)+p_i(a)p\right).
\]
Thus \(R_{a,i}\) is an affine exact symplectomorphism from the incoming
vertical port to the outgoing horizontal port \((0,1)\times J_i(a)\).

Let \(\Sigma_a\) be the graph completion of the open square after cutting
along every vertical seam \(q=s_i(a)\) and every horizontal seam
\(p=t_i(a)\).  A seam has two oriented copies and a corner has the four
corresponding quadrant copies.  Both
\[
 \mathcal Q_a^-=igsqcup_i I_i(a)\times(0,1),
 \qquad
 \mathcal Q_a^+=\bigsqcup_i(0,1)\times J_i(a)
\]
embed as dense open subsets of \(\Sigma_a\).  The coordinate identity on the
open square has a unique continuous lift
\[
 C_a:\overline{\mathcal Q_a^+}^{\,g}\longrightarrow
     \overline{\mathcal Q_a^-}^{\,g}
\]
which sends an oriented horizontal-side germ to the vertical-port germ
selected by its \(q\)-coordinate.  This is the recutting operation; it keeps
physical coordinates and changes only the cut chart.

\begin{lemma}[Symplectic recutting]
\label{lem:r5-a1-recut}
The map \(C_a\) is a bijection of graph completions, is smooth on every
corner stratum, preserves \(dq\wedge dp\), and depends \(C^3\) on \(a\) in
the moving-cut charts.  Consequently
\[
 F_a=C_a\circ R_a:\Sigma_a\longrightarrow\Sigma_a
\]
is a well-defined invertible piecewise exact symplectic return map.  Its
symbolic itinerary is the full four-symbol shift: the next symbol is selected
from the new \(q\)-coordinate and is not forced to equal the preceding port.
\end{lemma}

\begin{proof}
Refine the square into the sixteen rectangles
\(I_j(a)\times J_i(a)\).  On each closed quadrant copy the recut map is the
coordinate identity, hence is symplectic and bijective.  Two quadrant copies
are glued only when their oriented germs agree; the coordinate identity
preserves that relation, so the cell maps paste to a homeomorphism of the
graph completion.  The inverse is the same coordinate identity read in the
opposite cut atlas.  Moving a seam is conjugated to a fixed seam by the affine
normal coordinate
\[
 u=\frac{q-s_{j-1}(a)}{p_j(a)},
 \qquad
 v=\frac{p-t_{i-1}(a)}{p_i(a)}.
\]
In these charts all stratum maps are independent of \(a\), which proves the
claimed parameter regularity.  Finally \(R_{a,i}\) maps its branch onto the
whole horizontal port \(J_i(a)\); recutting then chooses any next vertical
port according to the unrestricted first coordinate.  Hence no branch label
is frozen.
\end{proof}

\subsection{An autonomous Hamiltonian suspension}

Let
\[
 \mathcal M_a=
 \bigl(\Sigma_a\times[0,1]_\tau\times\mathbb R_E\bigr)
 \Big/igl((x,1,E)\sim(F_ax,0,E)\bigr).
\]
On each stratum put
\[
 \Omega_a=dq\wedge dp+dE\wedge d\tau,
 \qquad H_a(q,p,\tau,E)=E.
\]
Because \(F_a\) is symplectic, both tensors descend through the gluing.
The Hamiltonian vector field is \(X_{H_a}=\partial_\tau\).

\begin{theorem}[Autonomous recut suspension]
\label{thm:r5-a1-suspension}
The quotient \(\mathcal M_a\) is a Hausdorff symplectic manifold with corners,
its corner smoothing may be chosen uniformly for \(a\) in every compact
parameter interval, and the section \(\Sigma_a\times\{0\}\times\{E_0\}\)
has first-return time one and first-return map exactly \(F_a\).
Moreover the affine branch correspondence is realized inside disjoint
Darboux port collars by a compactly supported time-dependent Hamiltonian,
so the suspension is symplectomorphic near the section to an autonomous
Hamiltonian impact network with the declared incoming and outgoing ports.
\end{theorem}

\begin{proof}
The equivalence relation is generated by a homeomorphism of the compact cut
section, so its graph is closed; the quotient is therefore Hausdorff.  Each
corner chart is a finite orthant crossed with \((0,1)\times\mathbb R\), and
the gluing is a smooth stratum-preserving diffeomorphism.  The standard
rounding function on an orthant can be selected in the fixed moving-cut
coordinates of Lemma~\ref{lem:r5-a1-recut}, giving a uniform smoothing.
Since \(dE\wedge d\tau\) and \(dq\wedge dp\) agree at the two boundary
parametrizations, \(\Omega_a\) descends and is nondegenerate.  The flow of
\(E\) moves only \(\tau\); at time one the gluing applies \(F_a\), proving the
return statement.

For completeness, the local branch realization is explicit.  Write
\(\ell_i=\log p_i(a)\).  The linear part of \(R_{a,i}\) is the time-one map of
\[
 K_{i}^{\rm lin}(q,p)=-\ell_i qp.
\]
A translation is the time-one map of a linear Hamiltonian.  Concatenating
these two Hamiltonians gives an exact isotopy from the identity to
\(R_{a,i}\).  Embed the four source and target rectangles in pairwise disjoint
Darboux collars and multiply the generators by cutoff functions equal to one
on the port cores and flat at the collar boundaries.  The cutoff error is
removed by the relative Moser equation
\(\iota_{Y_s}\omega=-\beta_s\), where \(\beta_s\) is the exact difference of
the two branch primitives and vanishes on the core and collar boundary.
The solution is compactly supported and preserves both prescribed regions.
The four corrections have disjoint support and therefore sum.  The recut
collar is the coordinate identity, so it introduces neither flux nor an
additional impact.  Suspending this compactly supported isotopy gives the
announced impact-network model.
\end{proof}

\subsection{A global port-work form}

Fix \(0<\delta<1/4\), choose
\(\rho\in C_c^\infty((\delta,2\delta))\) with
\(\int_0^1\rho(\tau)d\tau=1\), and choose smooth port functions
\(\chi_i\) which equal one on the core of incoming port \(i\), have disjoint
supports in its Darboux collar, and vanish near every seam.  For prescribed
port works \(\epsilon_i\), define on the suspension chart
\[
 \alpha_a=\sum_{i=1}^4\epsilon_i\chi_i(x)\rho(\tau)d\tau.
\]
It vanishes in neighborhoods of \(\tau=0,1\), hence its boundary traces agree
without imposing \(\epsilon_i=\epsilon_j\).

\begin{proposition}[Descending work form and section cocycle]
\label{prop:r5-a1-work}
The form \(\alpha_a\) descends to a globally defined smooth one-form on
\(\mathcal M_a\).  If a return segment enters through the core of port \(i\),
then
\[
 \int_{0}^{1}\alpha_a(X_{H_a})\,d\tau=\epsilon_i.
\]
For every regular orbit the accumulated mechanical work is therefore the
additive section cocycle
\[
 W_n(x)=\sum_{k=0}^{n-1}\epsilon_{\iota(F_a^kx)}.
\]
The construction remains valid when successive port signs differ.
\end{proposition}

\begin{proof}
Descent follows because \(\alpha_a\) is identically zero on two open collars
of the glued boundary.  During the support of \(\rho\), the Hamiltonian
trajectory remains in the same incoming port core and \(\chi_i=1\); all
other \(\chi_j\) vanish.  Integration gives \(\epsilon_i\).  Additivity under
concatenation gives the last formula.
\end{proof}

\subsection{Typed transfer and response operators}

Let \(\Sigma_4^+=\{1,2,3,4\}^{\mathbb N}\) with metric
\(d_\vartheta(x,y)=\vartheta^{n(x,y)}\), and let
\(\mathscr B=C^\eta(\Sigma_4^+)\).  For a normalized branch potential
\(\phi_a(i,x)\), define
\[
 (L_af)(x)=\sum_{i=1}^4e^{\phi_a(i,x)}f(ix),
 \qquad L_a\mathbf1=\mathbf1.
\]
For the benchmark Bernoulli family
\(e^{\phi_a(i,x)}=p_i(a)\).  Denote its invariant probability by \(\nu_a\)
and set
\[
 \Pi_af=\nu_a(f)\mathbf1,
 \qquad Q_a=I-\Pi_a,
 \qquad S_a=(I-L_a)^{-1}Q_a=\sum_{n\ge0}L_a^nQ_a.
\]
The series converges in \(\mathscr B\), because
\(\|L_a^nQ_a\|_{C^\eta\to C^\eta}\le C\vartheta^{\eta n}\).

\begin{theorem}[All-order typed response]
\label{thm:r5-a1-response}
Let \(K\) be a compact parameter interval on which
\(p_i(a)\ge p_*>0\) and \(a\mapsto\phi_a\) is \(C^{m+1}\) in
\(C^\eta\).  Then \(L_a,\Pi_a,Q_a,S_a\) are respectively bounded maps
\[
 L_a,\Pi_a,Q_a,S_a:\mathscr B\longrightarrow\mathscr B,
\]
are \(C^m\) in operator norm, and for every \(G_a\in C^m(K;\mathscr B)\)
all derivatives of \(\nu_a(G_a)\) up to order \(m\) are finite sums of typed
words made from
\[
 D^rL_a,\quad D^rG_a,\quad S_a,\quad\Pi_a.
\]
For example,
\[
 \frac d{da}\nu_a(G_a)
 =\nu_a(DG_a)+\nu_a\bigl((DL_a)S_aQ_aG_a\bigr).
\]
Under the common-root coupling of the branch variables, the same derivatives
are the limits of centered finite differences.  In physical cut coordinates
the derivative contains, in addition, the moving-seam saltation trace
\[
 \sum_i s_i'(a)\,[G_a]_{s_i(a)}\,\mathfrak j_i,
\]
where \([G_a]_{s_i(a)}\) is the oriented jump and \(\mathfrak j_i\) is the
invariant seam flux.  No undefined operator occurs in the formula.
\end{theorem}

\begin{proof}
The spectral gap follows directly from contraction of variations on the
one-sided full shift.  The Neumann series defining \(S_a\) is therefore
normally convergent, uniformly on \(K\).  Differentiate
\(L_a\mathbf1=\mathbf1\), \(\nu_aL_a=\nu_a\), and
\(\nu_a(\mathbf1)=1\).  On the centered range this gives
\(D\nu_a=\nu_a(DL_a)S_aQ_a\), which yields the displayed identity.
Repeated differentiation and the resolvent identity
\[
 DS_a=S_a(DL_a-D\Pi_a)S_a-D\Pi_a
\]
produce only the listed typed words; normal convergence justifies every
interchange.  The common-root construction uses
\(i_a(u)=i\) when \(u\in(s_{i-1}(a),s_i(a))\).  Away from moving seams the
sample is locally constant in \(a\); crossing a seam contributes its oriented
jump times the crossing speed.  Integrating in \(u\) gives the saltation term,
and dominated convergence plus the spectral gap gives the all-order limit.
\end{proof}

\subsection{Preparation and conditional calibration}

For branch counts \(N_i(n)\), multinomial coefficient extraction gives the
process-level Gibbs conditioning principle on every interior frequency vector.
The canonical field \(\theta\) is the gradient of the multinomial rate and the
canonical path likelihood is
\[
 \frac{d\mathbb P_\theta}{d\mathbb P_0}
 =\exp\{\theta W_n-n[Q(\theta)-Q(0)]\}.
\]

\begin{theorem}[Mechanical--valuation compatibility]
\label{thm:r5-a1-calibration}
A smooth, cash-additive, law-invariant and strongly time-consistent certainty
equivalent generated by one utility is either conditional expectation or an
entropic certainty equivalent with one constant coefficient \(\vartheta\).
The mechanical preparation law selects \(\vartheta=\theta\) precisely when
its likelihood cocycle is required to equal the canonical mechanical cocycle
above in the same work units.  Without that compatibility axiom mechanics does
not identify an unrelated preference coefficient.
\end{theorem}

\begin{proof}
Cash additivity gives the translation equation for the utility; differentiating
twice shows that its absolute risk aversion is constant, hence the utility is
affine or exponential.  The conditional tower makes the exponential
coefficient time independent.  Equality of the two normalized likelihoods on
the nonconstant cocycle \(W_n\) forces
\((\vartheta-\theta)W_n\) to be constant, hence \(\vartheta=\theta\).
\end{proof}

\begin{theorem}[Round-five A1 closure]
\label{thm:r5-a1-main}
The four-port family has a well-defined recut symplectic return map, an
autonomous Hamiltonian suspension, a global collar-supported mechanical work
form, a common symbolic preparation law, and a fully typed all-order response
calculus.  The valuation coefficient is identified with the preparation field
only under the explicit canonical-cocycle compatibility theorem.
\end{theorem}

\begin{proof}
Combine Lemma~\ref{lem:r5-a1-recut},
Theorem~\ref{thm:r5-a1-suspension},
Proposition~\ref{prop:r5-a1-work}, and
Theorems~\ref{thm:r5-a1-response} and
\ref{thm:r5-a1-calibration}.
\end{proof}
